\section{The components on local operators from the defect flow}

In this section $O=\int\omega\,X_1\cdots X_n$ is a monomial local functional, gauge-variant ones included, $\delta$ a square-zero parameter of degree $-|O|$, $O(\varepsilon)=O-\partial_\delta I^{CT}_{\mathfrak{psl}(4|4)}[\varepsilon](O)$ (Definition \ref{def:c8-renorm} with $\psi=0$), and $\hbar F_1$ a solution of \eqref{eq:c12-eqs} through arity one along the generators considered. As in Definition \ref{def:c11-defect} (iii), $\mathrm{ct}_m(M)$ and $\mathrm{lc}(M)$ are the coefficients of $\varepsilon^{-m}$ and $\log\varepsilon$ in $\partial_\delta I^{CT}[\varepsilon](M)$, and $\mathrm{ct}^{(1)}_{x,m}(M)$, $\mathrm{lc}^{(1)}_x(M)$ those of the arity-one counterterm of $M$ along $x$. For a two-body kernel $\kappa$ we write
\begin{equation}
\sum_{a<b}\pm\kappa(X_a,X_b):=\sum_{1\le a<b\le n}(-1)^{|X_b|(|X_{a+1}|+\cdots+|X_{b-1}|)}\int\omega\,X_1\cdots X_{a-1}\,\kappa(X_a,X_b)\,X_{a+1}\cdots\widehat{X_b}\cdots X_n,
\label{eq:c12-pairsum}
\end{equation}
the sum over the unordered pairs with the Koszul signs of \eqref{eq:c10-decomp}.

\begin{definition}[Defect cycles of a pair and two-body kernels]
\label{def:c12-kernel}
For a generator $x$, two letters $X_a$, $X_b$ of any slots and $0<\varepsilon<T\le\infty$, the defect cycles of the pair along $x$ and their sum $\mathcal D_x(X_a,X_b;\varepsilon,T)$, amputated with the outputs open and read with the sign of Definition \ref{def:c11-defect} (ii), are those of Definition \ref{def:c11-defect} (i), (ii): chains from $X_a$ to $X_b$ with vertices of $I$, exactly one vertex given by a term of $S^{(1)}(x)$ in a placement of Table \ref{tab:c12-place}, edges $P(\varepsilon,T)$ and exactly one special edge. The two-body kernel of the pair along $x$ is
\begin{equation}
\Upsilon_x(X_a,X_b):=[\varepsilon^0]\,\mathcal D_x(X_a,X_b;\varepsilon,T),
\label{eq:c12-kernel}
\end{equation}
the coefficient of $\varepsilon^0$ of the expansion of Lemma \ref{lem:c12-special}. In arity zero, $\Upsilon_0(X_a,X_b)$ is the coefficient of $\varepsilon^0$ of the chains from $X_a$ to $X_b$ with vertices of $I$, edges $P(\varepsilon,T)$ and exactly one special edge. For even $x$ acting on $A|_V$ the defect cycles with the vertex on a line are the special-edge composite cycles of Definition \ref{def:c11-sp}, and their coefficient of $\varepsilon^0$ is the composite pair kernel, written $\mathcal C_x(X_a,X_b)$.
\end{definition}

\begin{lemma}[The defect vertex of an operator]
\label{lem:c12-dv}
Modulo $\hbar^2$ and constants, the defect vertex $\mathfrak d_O:=(d_{CE}+Q)(J+\delta O(\varepsilon))+\tfrac12\{J+\delta O(\varepsilon),J+\delta O(\varepsilon)\}_\varepsilon+\hbar\Delta_\varepsilon(J+\delta O(\varepsilon))$ has the parts: (a) the defect vertex of the action, independent of $\delta$; (b) in arity zero, $\delta$ times
\[
sO+\{I,O\}^\sigma_\varepsilon+\hbar\Delta_\varepsilon O-s\,\partial_\delta I^{CT}[\varepsilon](O)-\{I,\partial_\delta I^{CT}[\varepsilon](O)\}^\sigma_\varepsilon-\{I^{CT}[\varepsilon]+\hbar F_1^{(0)},O\}_\varepsilon;
\]
(c) in arity one along $x$, $\delta$ times
\[
\pm\{S^{(1)}(x),O\}\pm\{S^{(1)}(x),O\}^\sigma_\varepsilon-s\,\mathrm{ct}^{(1)}_x(O)[\varepsilon]-\{I,\mathrm{ct}^{(1)}_x(O)[\varepsilon]\}^\sigma_\varepsilon\mp\{S^{(1)}(x),\partial_\delta I^{CT}[\varepsilon](O)\}_\varepsilon-\big\{I^{CT,(1)}(x)+\hbar F_1^{(1)}(x),O\big\}_\varepsilon,
\]
with the signs of the cochain bracket of Definition \ref{def:c2-dce} on an arity-zero argument (Theorem \ref{thm:c8-L} (iv)), all $+$ for $|O|=0$, and $\mathrm{ct}^{(1)}_x(O)[\varepsilon]=\sum_{m\ge1}\varepsilon^{-m}\mathrm{ct}^{(1)}_{x,m}(O)+\log\varepsilon\,\mathrm{lc}^{(1)}_x(O)$.
\end{lemma}

\begin{proof}
Expand in $\delta$, in the arity and in $\hbar$, with $I_{\mathfrak{psl}(4|4)}=I+S^{(1)}+S^{(2)}$. The classical master equation \eqref{eq:c2-mc} removes the brackets without special edge in (a); $d_{CE}$ vanishes on arity zero, so the parts of arity zero and one of $(d_{CE}+Q)(\delta O)+\{I_{\mathfrak{psl}(4|4)},\delta O\}$ are $\delta\,sO$ and $\pm\delta\{S^{(1)}(x),O\}$. $\hbar\Delta_\varepsilon$ of the arity-one counterterms and of $\hbar F_1$, and the brackets of two counterterms, are of order $\hbar^2$; $\hbar\Delta_\varepsilon S^{(1)}(x)$ is a constant (Lemma \ref{lem:c12-defect}).
\end{proof}

\begin{lemma}[Bridges and their counterterm partners]
\label{lem:c12-bridge}
Let $G$ be a subgraph at a generic momentum $p$, joined to the rest of a graph by a special edge on a tree, with $G=\sum_{m\ge1}\varepsilon^{-m}a_m+a_0+o(1)$ and singular part $\operatorname{Sing}G=\sum_{m\ge1}\varepsilon^{-m}a_m$ (plus a multiple of $\log\varepsilon$). The bare bridge has the finite limit $\sum_{m\ge1}\frac{(-1)^m}{m!}|p|^{2m}a_m$. Its counterterm partner, the special-edge part of the counterterm bracket across the same edge, is $-(e^{-\varepsilon|p|^2}-1)\operatorname{Sing}G$, whose coefficient of $\varepsilon^0$ is the negative of that limit; for $G$ a cycle through two letters of $M$ with $S^{(1)}(x)$ on it, or a cycle through two letters of $M$ in arity zero, these coefficients are the finite special vertex of Definition \ref{def:c11-defect} (iii) and the finite vertex \eqref{eq:c10-finitesp}. The sum of the bridge and its partner is
\begin{equation}
(e^{-\varepsilon|p|^2}-1)\big(G-\operatorname{Sing}G\big)\longrightarrow0 .
\label{eq:c12-bridge}
\end{equation}
\end{lemma}

\begin{proof}
$e^{-\varepsilon|p|^2}-1=\sum_{j\ge1}(-\varepsilon|p|^2)^j/j!$ at a generic momentum; multiplying $\varepsilon^{-m}a_m$ by it gives the poles $j<m$, the finite term $j=m$ and terms that vanish for $j>m$; a term $\log\varepsilon$ gives $(e^{-\varepsilon|p|^2}-1)\log\varepsilon\to0$. The counterterm vertex across the edge carries $-\operatorname{Sing}G$, so the partner is $-(e^{-\varepsilon|p|^2}-1)\operatorname{Sing}G$, and in the sum the convergent graph $G-\operatorname{Sing}G$ is multiplied by a factor tending to $0$. The coefficients $-\frac{(-1)^m}{m!}\{f,g_m\}_{H^m}$ of the brackets defining Definition \ref{def:c11-defect} (iii) and \eqref{eq:c10-finitesp} are the $\varepsilon^0$ coefficients of the partners with $f=I$, $g_m=\mathrm{ct}^{(1)}_{x,m}$ or $\mathrm{ct}_m$, and $f=S^{(1)}(x)$, $g_m=\mathrm{ct}_m$.
\end{proof}

\begin{theorem}[Every component from the defect flow]
\label{thm:c12-master}
Let $x$ be a generator, $O=\int\omega\,X_1\cdots X_n$ of degree $0$, $\psi=0$, and $\hbar F_1$ a solution of \eqref{eq:c12-eqs} in arity zero and in arity one along $x$. Take $T=\infty$ for $x\neq K_b$ and any $T\in(0,\infty)$ for $x=K_b$, and write the left side of \eqref{eq:c8-rep} in arity one along $x$ as $\big(L(O)^{(1)}(x)\big)[T]+\big(L(O)^{(0)}\big)[T]^{(1)}(x)$, with $L(O)^{(1)}(x)=\{S^{(1)}(x),O\}+\hbar L_1(O)^{(1)}(x)$ and $L(O)^{(0)}=sO+\hbar L_1(O)^{(0)}$ (Theorem \ref{thm:c8-L} (iv)).
\begin{itemize}
\item (i) The $\delta$-linear graphs of arity one along $x$ and of order $\hbar$ of the flow of $\mathfrak d_O$ form the classes: (F) one-loop graphs with the vertex $\pm\delta\{S^{(1)}(x),O\}$, or with the vertex $\delta\,sO$ and a vertex $S^{(1)}(x)$; (C2) the defect cycles of the pairs of letters of $O$ with trees of $I$ on the outputs; (C0) cycles through two letters of $O$ with one special edge or one contraction through $K_\varepsilon$, the vertex $S^{(1)}(x)$ on a tree; (Br) graphs with a special edge on a tree; (Ac) graphs whose cycle passes through vertices of $J$ alone; (Ct) trees with one vertex of order $\hbar$ of (b) or (c) of Lemma \ref{lem:c12-dv}.
\item (ii) As $\varepsilon\to0$: (F) with the counterterm trees of $\{S^{(1)}(x),O\}$ and of $sO$ gives the one-loop parts of $(\{S^{(1)}(x),O\})[T]$ and of $(sO)[T]^{(1)}(x)$; (C0) with the arity-zero vertices of (Ct) tends to $\big(\hbar L_1(O)^{(0)}\big)[T]^{(1)}(x)$; (Br) with its counterterm partners in (Ct) tends to $0$ by Lemma \ref{lem:c12-bridge}, the partners being used up there; (Ac) tends to $0$ by the arity-zero and arity-one equations.
\item (iii) (Master formula.)
\begin{equation}
\hbar L_1(O)^{(1)}(x)=\sum_{a<b}\pm\Upsilon_x(X_a,X_b)-\big\{\hbar F_1^{(1)}(x),O\big\},\qquad \hbar L_1(O)^{(0)}=\sum_{a<b}\pm\Upsilon_0(X_a,X_b)-\big\{\hbar F_1^{(0)},O\big\}.
\label{eq:c12-master}
\end{equation}
\item (iv) The defect cycles contain no $\log\varepsilon$; the coefficient of $\log\varepsilon$ cancels by \eqref{eq:c11-logcov}, and the coefficients of $\varepsilon^{-m}$ by \eqref{eq:c11-pole} with the shift terms $j\ge1$ removed, these being the poles of the bridge partners grouped in (ii); $I^{CT,(1)}(x)$ has no pole along every generator.
\item (v) Along $K_b$ the right side of \eqref{eq:c12-master} is the same local functional at every finite $T$; along the other generators at every $T\in(0,\infty]$.
\end{itemize}
\end{theorem}

\begin{proof}
By Lemma \ref{lem:c8-transfer}, $(d_{CE}+Q+\hbar\Delta_T)e^{\hbar\partial_{P(\varepsilon,T)}}e^{(J+\delta O(\varepsilon))/\hbar}=e^{\hbar\partial_{P(\varepsilon,T)}}\big(\hbar^{-1}\mathfrak d_O\,e^{(J+\delta O(\varepsilon))/\hbar}\big)$, so the left side of \eqref{eq:c8-rep} at scale $T$ is the $\delta$-linear part of the $\varepsilon\to0$ limit of the connected graphs with one vertex $\mathfrak d_O$, the others from $J+\delta O(\varepsilon)$ and edges $P(\varepsilon,T)$ (Lemma \ref{lem:c8-flow}, Proposition \ref{prop:c4-defectflow}), and by Definition \ref{def:c8-rep} it equals $L(O)[T]$.

(i) At order $\hbar$ a connected graph has one cycle and vertices of order $\hbar^0$, or one vertex of order $\hbar$ and no cycle (Proposition \ref{prop:c3-graphs}); one vertex carries $\delta$ and one carries $x$. A cycle passes through $O$ along exactly two letters or avoids it (Theorem \ref{thm:c10-decomp} (i)); the positions of the special edge and of $S^{(1)}(x)$ relative to the cycle sort the graph into (C2), (C0), (Br) or (Ac), the vertices $\pm\delta\{S^{(1)}(x),O\}$ and $\delta\,sO$ into (F), and the trees with a vertex of order $\hbar$ are (Ct).

(ii) (F): graph by graph the one-loop graphs of the observables $\{S^{(1)}(x),O\}$ and $sO$ at scale $T$; the differences of their counterterm trees with the vertices of (b), (c) are local and singular and enter (iv). (C0): Theorem \ref{thm:c10-brst} (proof), with a tree carrying $S^{(1)}(x)$ attached at generic momenta, and the arity-zero part of the present argument. (Br): Lemma \ref{lem:c12-bridge}; the subgraph is a cycle through two letters of $O$ with $S^{(1)}(x)$ on it, a cycle through two letters in arity zero, or a cycle of $J$ alone such as the self-energy with the pole of $I^{CT}$. Along $K_b$ at finite $T$ the subgraphs carry at fixed $\varepsilon$ the term $\tfrac3{8\pi^2}\hbar\log T$ times a tree with the same factor (Theorem \ref{thm:c8-K}), which also tends to $0$. (Ac): a cycle of $J$ alone, with its counterterm and finite vertices, is a graph of the flow of the defect of the action attached to $O$ by a tree; its limit is the tree-level flow of $\hbar\,\mathrm{Ob}-s(\hbar F_1)$ (Theorem \ref{thm:c8-obstruction} (ii)), zero in arity zero and in arity one along $x$.

(iii) What remains is (C2) and the arity-one vertices of (Ct) without their bridge partners: $-s\,\mathrm{ct}^{(1)}_x(O)[\varepsilon]$, $\mp\{S^{(1)}(x),\partial_\delta I^{CT}[\varepsilon](O)\}$, $-\{I^{CT,(1)}(x),O\}$, $-\{\hbar F_1^{(1)}(x),O\}_\varepsilon$, with the counterterm trees $-\partial_\delta I^{CT}[\varepsilon](\{S^{(1)}(x),O\})$ and $-\mathrm{ct}^{(1)}_x(sO)$ moved from the left side. Each defect cycle is $\sum_{m\ge1}\varepsilon^{-m}a_m+a_0+o(1)$ with local coefficients and no $\log\varepsilon$ (Lemma \ref{lem:c12-special}, Theorem \ref{thm:c11-oddsimplex}); the trees on its outputs converge to trees with edges $P(0,T)$. $-\{\hbar F_1^{(1)}(x),O\}_\varepsilon$ gives $-\{\hbar F_1^{(1)}(x),O\}$, its special-edge part tending to $0$. After the cancellations of (iv) the sum converges to the tree flow of the right side of \eqref{eq:c12-master}, and Theorem \ref{thm:c8-L} (i) identifies the local functionals; the sum over the pairs with the Koszul signs is Theorem \ref{thm:c10-decomp} (iii). In arity zero the same argument with the vertices (b) gives the second formula; the finite part of $\{I^{CT}[\varepsilon],O\}^\sigma_\varepsilon$ is the partner of the bridges through the self-energy cycles and cancels with them.

(iv) The left side has a limit, so the coefficients of $\log\varepsilon$ and of $\varepsilon^{-m}$ of the collected vertices vanish, which with the injectivity of the tree flow gives \eqref{eq:c11-logcov} and \eqref{eq:c11-pole} in the stated form. A coefficient of $\varepsilon^{-m}$ of $I^{CT,(1)}(x)$ has degree $-1$ and total Weyl weight $4-\tfrac12\lambda_x-2m$ with position degree at most $p_x$; a colour singlet of degree $-1$ has total weight at least $\tfrac72-p_x$, while for $m\ge1$ the total weight is at most $2,3,2,2,\tfrac52,\tfrac32,1$ along $\Delta,P_b,M_\omega,R_r,\Sigma,\Sigma^*,K_b$, below $\tfrac52,\tfrac72,\tfrac52,\tfrac72,\tfrac72,\tfrac52,\tfrac32$ (Proposition \ref{prop:c12-poles}).

(v) Lemma \ref{lem:c12-special} and Theorem \ref{thm:c11-oddsimplex} (iii); $\hbar F_1^{(1)}(K_b)$ is independent of $T$ (Proposition \ref{prop:c12-T}); along the other generators Proposition \ref{prop:c8-Tsafe} and Theorem \ref{thm:c8-L} (ii).
\end{proof}

\begin{remark}[The finite special vertices in Chapters 10 and 11]
\label{rem:c12-finitesp}
Lemma \ref{lem:c12-bridge} pairs the finite special vertex of Definition \ref{def:c11-defect} (iii) and the finite vertex \eqref{eq:c10-finitesp} with the bare bridges: both are the coefficients of $\varepsilon^0$ of the counterterm partners of the bridges, and they cancel against the bridges. The two-body kernels of the components are therefore the coefficients of $\varepsilon^0$ of the defect cycles alone, \eqref{eq:c12-kernel}, and the pole identity holds in the form of Theorem \ref{thm:c12-master} (iv). On the leg sets with superficial degree $D<2$ the defect cycles have no poles and the finite special vertex vanishes (Proposition \ref{prop:c11-oddbound}), so there the kernels of Definition \ref{def:c11-defect} (iv) and of \eqref{eq:c12-kernel} coincide; the composite pair kernels of Definition \ref{def:c11-sp} are coefficients of $\varepsilon^0$ of cycles and are the same in both readings. On the computed leg sets with $D\ge2$ the pair form of Remark \ref{rem:c12-pair} holds with \eqref{eq:c12-kernel}, in both evaluations of Remark \ref{rem:c12-method}.
\end{remark}

Every number of a component is a coefficient of $\varepsilon^0$ of a defect cycle or a coefficient of the action-side solution: the master formula uses Lemma \ref{lem:c8-transfer}, Definition \ref{def:c8-rep}, the injectivity of the tree flow, the classical master equation, the equations \eqref{eq:c12-eqs} through arity one along $x$, the simplex form of the defect cycles and generic momenta. The identities below are proved from it.

\begin{theorem}[Ward identity of the flow along an even generator]
\label{thm:c12-ward}
Let $x$ be even, $J$ even and at least cubic modulo $\hbar$, $P=P(\varepsilon,T)$ and $W=W(P,J)$. With $\{S^{(1)}(x),-\}=x.$ on functionals (Definition \ref{def:c7-vdot}),
\begin{equation}
\{S^{(1)}(x),W\}=\partial_\delta W\big(P,J+\delta\{S^{(1)}(x),J\}\big)-\hbar\,\partial_{\mathcal V_x}W-\tfrac12\big(\partial_{\mathcal V_x}(W^2)-2W\,\partial_{\mathcal V_x}W\big),
\label{eq:c12-ward}
\end{equation}
the last two terms being the graphs of $W$ with one edge $P$ replaced by the variation kernel $\mathcal V_x$. In particular:
\begin{itemize}
\item (i) $\{S^{(1)}(\Delta),W\}=\partial_\delta W(P,J+\delta\{S^{(1)}(\Delta),J\})+4(\varepsilon\partial_\varepsilon+T\partial_T)W$, the derivatives acting on $P$, which is \eqref{eq:c9-ward};
\item (ii) for $x\in\{P_b,M_\omega,R_r\}$, $\{S^{(1)}(x),W\}=\partial_\delta W(P,J+\delta\{S^{(1)}(x),J\})$ at every $\varepsilon$ and $T$, which is \eqref{eq:c7-flowcov}.
\end{itemize}
\end{theorem}

\begin{proof}
By Lemma \ref{lem:c9-ops}, $[\partial_\phi,x.]=\partial_{(x.\otimes1+1\otimes x.)\phi}$ for even $x.$, so $x.\,\partial_P=\partial_P\,x.-\partial_{\mathcal V_x}$; the operators $\partial_\phi$ commute, hence $x.\,e^{\hbar\partial_P}=e^{\hbar\partial_P}(x.-\hbar\partial_{\mathcal V_x})$. Apply this to $e^{J/\hbar}$: $x.\,e^{W/\hbar}=\hbar^{-1}(x.W)e^{W/\hbar}$, $e^{\hbar\partial_P}(x.e^{J/\hbar})=\hbar^{-1}\partial_\delta W(P,J+\delta\,x.J)e^{W/\hbar}$, and $\partial_{\mathcal V_x}e^{W/\hbar}=\big(\hbar^{-1}\partial_{\mathcal V_x}W+\tfrac12\hbar^{-2}(\partial_{\mathcal V_x}(W^2)-2W\partial_{\mathcal V_x}W)\big)e^{W/\hbar}$ (Definition \ref{def:c3-delta}); multiply by $\hbar e^{-W/\hbar}$. (i) and (ii) are Proposition \ref{prop:c12-composite} (ii) and (i).
\end{proof}

\begin{lemma}[The arity-one obstruction along an even generator]
\label{lem:c12-obeven}
For $x\in\{\Delta,P_b,M_\omega,R_r\}$, with $\hbar F_1^{(0)}$ a solution of the arity-zero equation invariant under $x$, $I^{CT,(1)}(x)$ is the singular part of the composite loops of $I$, and
\begin{equation}
\hbar\,\mathrm{Ob}^{(1)}(x)=\mathcal V^{\rm act}_x+(\text{composite part along }x),\qquad \mathcal V^{\rm act}_\Delta=4L^{(0)}=4\,s\mathcal N\ \text{units},\qquad \mathcal V^{\rm act}_{P_b}=\mathcal V^{\rm act}_{M_\omega}=\mathcal V^{\rm act}_{R_r}=0,
\label{eq:c12-obeven}
\end{equation}
the composite part along $x$ being the local functional whose tree flow is $\lim_{\varepsilon\to0}\big(Q(\text{composite loops of }I)-\text{tree flow of }s\,I^{CT,(1)}(x)\big)$; it is $0$ along $R_r$, so $\mathrm{Ob}^{(1)}(R_r)=0$ and $I^{CT,(1)}(R_r)=0$.
\end{lemma}

\begin{proof}
As Proposition \ref{prop:c9-structure}, with Theorem \ref{thm:c12-ward} in place of \eqref{eq:c9-D1a}, \eqref{eq:c9-D1b} and \eqref{eq:c9-ward}. By Theorem \ref{thm:c12-second} (ii) the linearised flow of $S^{(1)}(x)$ is the vertex, the vertex joined by one edge to a leg of $W$, and the composite closures; $Q$ of it is $\{S^{(1)}(x),W\}_\varepsilon$ minus the terms with $QW$, using $QP\star x.+P\star x.Q=E(H)x.$ and $\{S^{(1)}(x),F\}_t=(e^{-tH}x.)F$ (Lemma \ref{lem:c9-SDelta} (ii)). The special-edge part on legs tends to $0$, and by \eqref{eq:c12-ward} $\{S^{(1)}(x),W\}$ is the linearised flow of $x.J$ plus the variation term. For $\Delta$, $x.I=x.\hbar F_1^{(0)}=0$, $(v.-4\varepsilon\tfrac d{d\varepsilon})I^{CT}[\varepsilon]=-4L^{(0)}$ (proof of Corollary \ref{cor:c9-CS}), and $4\varepsilon\partial_\varepsilon W\to0$, which gives $\mathcal V^{\rm act}_\Delta$. For $P_b$, $M_\omega$, $R_r$, $x.I=0$, $x.I^{CT}[\varepsilon]=0$ by the covariance of the counterterms (Proposition \ref{prop:c10-map} (iv)), $x.\hbar F_1^{(0)}=0$, and the variation kernel vanishes. The terms with $QW$ tend to $0$ by the arity-zero equation, except $Q$ of the composite loops, $Q$ anticommuting with $\partial$ of the composite line (Proposition \ref{prop:c12-composite}). Along $R_r$ the composite line is absent (Table \ref{tab:c12-place}).
\end{proof}

\begin{theorem}[The conjugation evaluation along $\Delta$, $P_b$, $M_\omega$, $R_r$]
\label{thm:c12-conj}
Let $x\in\{\Delta,P_b,M_\omega,R_r\}$, $\hbar F_1^{(0)}$ invariant under $x$, and $O$ homogeneous with $sO=0$ and $\hbar L_1(O)^{(0)}=0$ for a finite part $\psi$ of the operators that is linear in $\omega$, translation invariant, weight preserving and covariant under $\mathfrak{so}(V)\oplus\mathfrak{so}(W_{\mathbb R}^\perp)$. The composite term of $O$ along $x$ is the local functional whose tree flow is the limit of $Q$ of the composite graphs of $O$ minus their counterterm trees, minus the tree insertion of the composite part into $O$ (Definition \ref{def:c10-hyp} along $\Delta$). Then at every $T\in(0,\infty]$
\begin{equation}
\hbar L_1(O)^{(1)}(x)=\mathcal V_x(O)+(\text{composite term of }O\text{ along }x)-\big\{\hbar F_1^{(1)}(x),O\big\},\qquad \mathcal V_\Delta(O)=4\,\mathrm{lc}(O),\quad \mathcal V_{P_b}=\mathcal V_{M_\omega}=\mathcal V_{R_r}=0 .
\label{eq:c12-conjev}
\end{equation}
\end{theorem}

\begin{proof}
At $T=\infty$, $d_\infty=d_{CE}+Q$, and the arity-one component along $x$ of $d_\infty O[\infty]$ is $Q$ of the arity-one part of $O[\infty]$, the limit of the $\delta$-part of the linearised flow at $J+\delta O(\varepsilon)$ of $S^{(1)}(x)-I^{CT,(1)}(x)-\hbar F_1^{(1)}(x)-\delta\,\mathrm{ct}^{(1)}_x(O)[\varepsilon]$. (1) By Theorem \ref{thm:c12-ward} the vertex $S^{(1)}(x)$ gives the linearised flow of $x.J$ plus the variation term; the $\delta$-part of $x.J$ is $x.O(\varepsilon)$. For $\Delta$, $(v.-4\varepsilon\tfrac d{d\varepsilon})O(\varepsilon)=(v.O)(\varepsilon)+4\,\mathrm{lc}(O)$ by \eqref{eq:c10-scaling}, and the $\delta$-independent part gives $4$ times the tree insertion of $L^{(0)}$ into $O[\infty]$; for $P_b$, $M_\omega$, $R_r$, $x.\partial_\delta I^{CT}[\varepsilon](O)=\partial_\delta I^{CT}[\varepsilon](x.O)$ (Proposition \ref{prop:c10-map} (iv)) and $x.\psi=\psi x.$. The terms with $QW$ vanish in the limit by the arity-zero equation, $sO=0$ and $L_1(O)^{(0)}=0$, except $Q$ of the composite graphs. (2) The vertex $-\hbar F_1^{(1)}(x)$ enters by tree insertion; by Lemma \ref{lem:c6-bvexact} with $sO=0$, $-Q$ of its tree insertion is minus the tree insertion of $s\hbar F_1^{(1)}(x)=\hbar\,\mathrm{Ob}^{(1)}(x)$, which by Lemma \ref{lem:c12-obeven} is $\mathcal V^{\rm act}_x$ plus the composite part, minus $\{\hbar F_1^{(1)}(x),O\}[\infty]$; the first cancels the $\delta$-independent term of (1) along $\Delta$. (3) The vertices $-I^{CT,(1)}(x)$ and $-\delta\,\mathrm{ct}^{(1)}_x(O)$ are the counterterm trees of the composite graphs. With $\{S^{(1)}(x),O\}=x.O$, Theorem \ref{thm:c8-L} (i) gives \eqref{eq:c12-conjev}; along $R_r$ the composite graphs vanish, and Theorem \ref{thm:c8-L} (ii) gives every finite $T$.
\end{proof}

\begin{theorem}[Dilatation, translations and rotations from the defect flow]
\label{thm:c12-D1}
Under the hypotheses of Theorem \ref{thm:c12-conj}, for every cocycle $O=\int\omega\,X_1\cdots X_n$:
\begin{equation}
\begin{aligned}
D_1(O)&=4\big(\mathrm{lc}(O)-\{\mathcal N,O\}\ \text{units}\big)+\sum_{a<b}\pm\mathcal C_\Delta(X_a,X_b)-\big\{\hbar F_1^{(1)}(\Delta)-4\,\mathcal N\ \text{units},O\big\},\\
P_1(O)&=\sum_{a<b}\pm\mathcal C_{P_b}(X_a,X_b)-\big\{\hbar F_1^{(1)}(P_b),O\big\},\qquad M_1(O)=\sum_{a<b}\pm\mathcal C_{M_\omega}(X_a,X_b)-\big\{\hbar F_1^{(1)}(M_\omega),O\big\},\\
R_1(O)&=-\big\{\hbar F_1^{(1)}(R_r),O\big\},
\end{aligned}
\label{eq:c12-D1}
\end{equation}
and, summed over the pairs of the cocycle,
\begin{equation}
\sum_{a<b}\pm\Upsilon_\Delta=\sum_{a<b}\pm\big(4\,\mathrm{lc}(X_aX_b)+\mathcal C_\Delta(X_a,X_b)\big),\quad \sum_{a<b}\pm\Upsilon_{P_b}=\sum_{a<b}\pm\mathcal C_{P_b},\quad \sum_{a<b}\pm\Upsilon_{M_\omega}=\sum_{a<b}\pm\mathcal C_{M_\omega},\quad \sum_{a<b}\pm\Upsilon_{R_r}=0 .
\label{eq:c12-pairsums}
\end{equation}
\end{theorem}

\begin{proof}
Theorems \ref{thm:c12-master} (iii) and \ref{thm:c12-conj} compute the same local representative, which is unique (Theorem \ref{thm:c8-L} (i)), and the one-body terms $-\{\hbar F_1^{(1)}(x),O\}$ are the same in both. By Theorem \ref{thm:c12-local} (iv), with Remark \ref{rem:c12-localM} for $M_\omega$, the composite term of a cocycle is $\sum_{a<b}\pm\mathcal C_x(X_a,X_b)$. For $\Delta$, $4\,\mathrm{lc}(O)=\sum_{a<b}\pm4\,\mathrm{lc}(X_aX_b)$ (Theorem \ref{thm:c10-decomp} (iii)), and the map $M\mapsto\mathrm{lc}(M)-\{\mathcal N,M\}$ units is that of Theorem \ref{thm:c10-brst}.
\end{proof}

The factor four and the logarithmic coefficient in \eqref{eq:c12-D1} are outputs of the theorem, an identity between two evaluations of one defect flow: \eqref{eq:c12-master}, which defines the numbers, and the evaluation through the regulator variation of Theorem \ref{thm:c12-ward}, a property of $P(\varepsilon,T)$ proved from its definition through \eqref{eq:c9-MP} and \eqref{eq:c10-scaling}, the latter following from $\sigma_\lambda P(\varepsilon,\infty)=P(\lambda^2\varepsilon,\infty)$ (Proposition \ref{prop:c6-scaling}).

\begin{remark}[The pair form]
\label{rem:c12-pair}
The pair form is the statement, for every pair of letters of any slots with any derivatives, that $\Upsilon_\Delta=4\,\mathrm{lc}+\mathcal C_\Delta$, $\Upsilon_{P_b}=\mathcal C_{P_b}$, $\Upsilon_{M_\omega}=\mathcal C_{M_\omega}$ and $\Upsilon_{R_r}=0$. The pair kernels are graded symmetric, and the difference $\Upsilon_x-\mathcal V_x-\mathcal C_x$ is a two-body kernel whose sum over the pairs of every cocycle vanishes by \eqref{eq:c12-pairsums}. On the pair $(A|_{W_{\mathbb R}^\perp},A|_{W_{\mathbb R}^\perp})$ the pair form follows from \eqref{eq:c12-pairsums}, the scalar bilinears being cocycles with this single pair: there $\Upsilon^0_\Delta=\tfrac{40}3O$ (traceless) and $-\tfrac{104}3O$ (trace) units, with $4\,\mathrm{lc}=16O$ and $-32O$ units and $\mathcal C^0_\Delta=-\tfrac83O$ units (Example \ref{ex:c11-bilinear}), and $\Upsilon_{P_b}=\mathcal C_{P_b}=-\tfrac13\partial_bO$ units (Corollary \ref{cor:c11-DPdefect}); along $R_r$ graded symmetry and $\mathfrak{so}(W_{\mathbb R}^\perp)$-covariance leave the form $a\big((R_r.A|_{W_{\mathbb R}^\perp})A|_{W_{\mathbb R}^\perp}+A|_{W_{\mathbb R}^\perp}(R_r.A|_{W_{\mathbb R}^\perp})\big)$, which the cocycle identity on the traceless bilinears sets to $0$. On every pair of the tables of Remark \ref{rem:c12-method} the pair form holds entry by entry in both evaluations. A proof pair by pair follows from Theorem \ref{thm:c12-conj} for the functional $\int\omega\,X_aX_b$, which is not a cocycle, with the terms in $s\int\omega\,X_aX_b$ of both evaluations, given by Theorem \ref{thm:c12-local} (vi) in the second and by the classes (F), (C0) in the first.
\end{remark}

\begin{definition}[Centred Killing data]
\label{def:c12-centred}
For a point $x_O$ and $z:=x-x_O$ the Killing data with explicit position decompose as
\[
\Delta:\ 2x=2x_O+2z,\qquad M_\omega:\ \omega x=\omega x_O+\omega z,\qquad \eta:\ \rho(x)\eta=\rho(x_O)\eta+\rho(z)\eta,\qquad K_b:\ v(x)=v(x_O)+\big(-2(b,x_O)z+\omega_{x_O}z\big)+v(z),
\]
with $\omega_{x_O}z:=2\big((x_O,z)b-(b,z)x_O\big)$. For $y$ with Killing data of positive degree, $\Upsilon^0_y(X_a,X_b)$ is the two-body kernel at the point $x_O$ of the pair computed with the Killing data centred at $x_O$; on plane waves, with the momentum $r$ injected at the vertex and absorbed by $\omega$, the monomials $z^\mu$, $z^\mu z^\nu$ act as $-i\partial_{r_\mu}$ and $-\partial_{r_\mu}\partial_{r_\nu}$ at $r=0$.
\end{definition}

\begin{theorem}[Position decompositions]
\label{thm:c12-pos}
\begin{itemize}
\item (i) For every pair of letters at the point $x_O$, graph by graph,
\begin{equation}
\begin{aligned}
&\Upsilon_\Delta=2x_O^\mu\,\Upsilon_{P_{e_\mu}}+\Upsilon^0_\Delta,\qquad \Upsilon_{M_\omega}=(\omega x_O)^\mu\,\Upsilon_{P_{e_\mu}}+\Upsilon^0_{M_\omega},\qquad \Upsilon_\eta=\sum_\mu x_O^\mu\,\Upsilon_{\rho(e_\mu)\eta}+\Upsilon^0_\eta,\\
&\Upsilon_{K_b}=v(x_O)^\mu\,\Upsilon_{P_{e_\mu}}-(b,x_O)\,\Upsilon^0_\Delta+\Upsilon^0_{M_{\omega_{x_O}}}+\Upsilon^0_{K_b},
\end{aligned}
\label{eq:c12-posdec}
\end{equation}
and the same decompositions hold for every coefficient $\mathrm{ct}^{(1)}_{x,m}$, $\mathrm{lc}^{(1)}_x$ and for $I^{CT,(1)}(x)$; $\Upsilon^0_\eta$ is computed with $S^{(1)}(\eta)$ at the spinor $\rho(z)\eta$ together with $s_0$, the term $Y_\eta$ of $\eta.$ belonging to it.
\item (ii) If $\hbar F_1$ solves the arity-two equation along $(y,P_b)$ for $y\in\{\Delta,M_\omega,\eta,K_b\}$, then, with $S^{(2)}(y,P_b)=0$, the position part along $b$ of $\hbar F_1^{(1)}(y)$ is
\begin{equation}
\pm\hbar F_1^{(1)}([y,P_b])\pm\big\{S^{(1)}(y),\hbar F_1^{(1)}(P_b)\big\}\pm s\big(\hbar F_1^{(2)}(y,P_b)\big)\mp\hbar\,\mathrm{Ob}^{(2)}(y,P_b),
\label{eq:c12-onebodypos}
\end{equation}
with $[\Delta,P_b]=-2P_b$, $[M_\omega,P_b]=-P_{\omega b}$, $[\eta,P_b]=-\rho(b)\eta$ and $[K_b,P_{b'}]=(b,b')\Delta+M_{bb'}$; in a scheme with $\hbar F_1^{(1)}(P_b)=0$ and an $s$-exact arity-two obstruction along $(y,P_b)$ the one-body term $-\{\hbar F_1^{(1)}(y),O\}$ decomposes as \eqref{eq:c12-posdec} up to $s$-exact brackets, zero on cohomology (Lemma \ref{lem:c10-deriv}).
\item (iii) If the local representatives exist through arity two along $(y,P_b)$ on a family stable under $\Lambda$ (Corollary \ref{cor:c8-homotopy}), then on its cocycles, with $\Lambda^{(1)}(x)=\Lambda^{(1)}_0(x)+\hbar\Lambda^{(1)}_1(x)$ as in Proposition \ref{prop:c11-relations},
\begin{equation}
\big[\Lambda^{(1)}_1(y),\Lambda^{(1)}_0(P_b)\big]+\big[\Lambda^{(1)}_0(y),\Lambda^{(1)}_1(P_b)\big]+\Lambda^{(1)}_1([y,P_b])\cong0
\label{eq:c12-poscoh}
\end{equation}
modulo $\Lambda^{(0)}$-exact terms: the position part of $D_1$ vanishes on cohomology when $P_1$ is translation covariant; the position part of $S_1$ along $\eta$ is, on cohomology, $Q_1$ along $\rho(b)\eta$ up to the commutator of $\{S^{(1)}(\eta),-\}$ with $P_1$; the position part of $K_1$ along $b'$ is, on cohomology, the component along $(b,b')\Delta+M_{bb'}$ up to the commutator of $\{S^{(1)}(K_b),-\}$ with $P_1$.
\end{itemize}
\end{theorem}

\begin{proof}
(i) $S^{(1)}(y)$ is linear in the Killing data: for even $y$ through $v.=v^\mu\partial_\mu+\tfrac12\omega_{\mu\nu}\Sigma^{\mu\nu}+\tfrac w4\operatorname{div}v$, whose terms are linear in $v$, $dv$ and $\operatorname{div}v$; for $\eta$ through \eqref{eq:c11-Seta}. Insert Definition \ref{def:c12-centred}: on plane waves $x^\mu=x_O^\mu+z^\mu$ acts as $x_O^\mu$ on the delta function of the total momentum absorbed by $\omega$ and as $-i\partial_{r_\mu}$ on the integrand (Proposition \ref{prop:c11-position} (i)); the constant part $v(x_O)$ of the field of $K_b$ is a translation, its linear part is $-(b,x_O)$ times the dilatation field $2z$ plus the rotation $\omega_{x_O}z$, with $\operatorname{div}v(x)=-8(b,x_O)-8(b,z)$, $\operatorname{div}(2z)=8$, $\operatorname{div}(\omega_{x_O}z)=0$, and the quadratic part $v(z)$ is $K_b$ centred at $x_O$. The decomposition commutes with the loop integral, the simplex integration, the special edge and the expansion in $\varepsilon$ (Theorem \ref{thm:c11-oddsimplex} (ii), (iii)), and the counterterms are built linearly from $S^{(1)}(y)$. (ii) The arity-two component of $s(\hbar F_1)=\hbar\,\mathrm{Ob}$ at $(y,P_b)$ has the terms $\pm\hbar F_1^{(1)}([y,P_b])$, $\pm\{S^{(1)}(y),\hbar F_1^{(1)}(P_b)\}\pm\{S^{(1)}(P_b),\hbar F_1^{(1)}(y)\}$ and $(Q+\{I,-\})\hbar F_1^{(2)}(y,P_b)$; $\{S^{(1)}(P_b),G\}$ for $G=\int(x^\mu g_\mu+g_0)$ is minus $\int b^\mu g_\mu$, the position part. (iii) Corollary \ref{cor:c8-homotopy} (ii) at $(y,P_b)$ with $S^{(2)}(y,P_b)=0$, whose order-$\hbar$ part is a graded commutator with $\Lambda^{(0)}$; Lemma \ref{lem:c11-canonical} identifies $[\mathcal M,\{S^{(1)}(P_b),-\}]$ with the position part of a map $\mathcal M$ polynomial in the position.
\end{proof}

\begin{corollary}[Special supercharges and special conformal fields]
\label{cor:c12-SK}
For $O=\int\omega\,X_1\cdots X_n$ at the point $x_O$ of its density, with the superscript $0$ denoting the parts centred at $x_O$,
\begin{equation}
\begin{aligned}
S_1(O)&=Q_1(O)\big|_{\rho(x_O)\eta}+S^0_1(O),\qquad S^0_1(O):=\sum_{a<b}\pm\Upsilon^0_\eta(X_a,X_b)-\big\{\hbar F_1^{(1)}(\eta)^0,O\big\},\\
K_1(O)&=v(x_O)^\mu P_1(O)\big|_{P_{e_\mu}}-(b,x_O)\,D^0_1(O)+M^0_1(O)\big|_{\omega_{x_O}}+K^0_1(O),
\end{aligned}
\label{eq:c12-SK}
\end{equation}
when the one-body terms decompose as in Theorem \ref{thm:c12-pos} (ii), and otherwise up to the bracket with the right side of \eqref{eq:c12-onebodypos} minus its first term. On the scalar bilinears $\Upsilon^0_\eta(A|_{W_{\mathbb R}^\perp},A|_{W_{\mathbb R}^\perp})$ vanishes on the colour singlet and $\{\int s_0,A|_{W_{\mathbb R}^\perp}\}=0$, so the two-body part of $S_1$ is that of $Q_1$ along $\rho(x_O)\eta$; $\Upsilon^0_{K_b}(A|_{W_{\mathbb R}^\perp},A|_{W_{\mathbb R}^\perp})$ and the rotation part vanish there (Proposition \ref{prop:c11-SKbil}).
\end{corollary}

\begin{proof}
Insert \eqref{eq:c12-posdec} into \eqref{eq:c12-master} for $y=\eta$, $K_b$, with the decompositions of the one-body terms; the components along $P_{e_\mu}$, $\rho(e_\mu)\eta$, $\Delta$ and $M_{\omega_{x_O}}$ are those of \eqref{eq:c12-master} along these generators.
\end{proof}

\begin{proposition}[The endpoint along $K_b$]
\label{prop:c12-KT}
Along $K_b$ at every finite $T$: every defect cycle has $T\partial_T=O(\varepsilon)$, so $\Upsilon_{K_b}$ is independent of $T$; the graphs of the classes (Br) and (Ac) carry at fixed $\varepsilon$ the term $\tfrac3{8\pi^2}\hbar\log T$ times a tree with the factor $e^{-\varepsilon|p|^2}-1$ at a generic momentum, or with its gradient $-2\varepsilon p\,e^{-\varepsilon|p|^2}$ along $(\eta_1,\eta_2)$, and tend to $0$; hence $K_1(O)$ is independent of $T$, and the observables keep the term $\tfrac3{8\pi^2}\hbar\log T$ of their order-$\hbar^0$ part generated by the flow (Theorem \ref{thm:c8-K}), whose image under $Q$ vanishes on gauge-invariant $O$.
\end{proposition}

\begin{proof}
Theorem \ref{thm:c12-master} (v). The coefficient of $\log T$ of a graph with $S^{(1)}(K_b)$ on its cycle is $\tfrac3{8\pi^2}\hbar$ times the tree obtained by cutting the insertion out of the cycle with the freed legs on the constant modes (Proposition \ref{prop:c8-Tsafe}); in (Br) and (Ac) this tree contains the special edge at a generic momentum. Theorem \ref{thm:c8-K}.
\end{proof}

\section{Locality of the composite term}

Throughout this section $x\in\{\Delta,P_b,M_\omega\}$, so that the composite line $\mathsf L_x:=(P\star x.\otimes1+1\otimes P\star x.)P$ is $Q$-closed (Proposition \ref{prop:c12-composite}), $T=\infty$ with generic momenta or any finite $T$, and $P=P(\varepsilon,T)$. The composite term enters the second evaluation (Theorem \ref{thm:c12-conj}) and the identities of Theorem \ref{thm:c12-D1}; the components themselves are given by Theorem \ref{thm:c12-master} along every generator. The theorem of this section completes the outline of Theorem \ref{thm:c11-local}.

\begin{definition}[Composite loops of the action]
\label{def:c12-Xi}
$\Xi_x(\varepsilon)$ is the sum of the amputated one-particle-irreducible composite loops of $I$ along $x$: cycles of vertices of $I$ and edges $P$ in which one edge is $\mathsf L_x$, all other legs open, so that $\tfrac\hbar2\partial_{\mathsf L_x}W(P,I)_{\rm tree}$ is the tree flow of $\Xi_x(\varepsilon)$. $\Xi^\sigma_x(\varepsilon)$ is minus the sum of the same loops with one edge of the cycle other than $\mathsf L_x$ replaced by the special edge, with the sign of Proposition \ref{prop:c11-master}. The renormalised composite functional is $\Xi^{\rm ren}_x:=\lim_{\varepsilon\to0}\big(\Xi_x(\varepsilon)-I^{CT,(1)}(x)\big)$, a functional of degree $-1$ at generic momenta, whose quadratic part along $\Delta$ is the renormalised composite two-point functional of Remark \ref{rem:c9-Pi}. The composite graphs of a local functional $O$ are of type I when the cycle passes through two letters of $O$ and of type II when it passes through vertices of $I$ alone; type II is the tree insertion of $\Xi_x(\varepsilon)$ into $O$, written $\mathrm{tr}(\Xi_x,O)$.
\end{definition}

\begin{lemma}[The image of the composite loops under $s$]
\label{lem:c12-Cid}
For every $\varepsilon>0$, as functionals at generic momenta, $s\,\Xi_x(\varepsilon)=\Xi^\sigma_x(\varepsilon)$. Consequently: (i) $\Xi^\sigma_x$ has no term in $\log\varepsilon$ and no poles, and $I^{CT,(1)}(x)=\log\varepsilon\cdot L^{(1)}(x)$ with $sL^{(1)}(x)=0$; (ii) the composite part along $x$ of Lemma \ref{lem:c12-obeven} is $s\,\Xi^{\rm ren}_x=[\varepsilon^0]\Xi^\sigma_x$, a local functional, although $\Xi^{\rm ren}_x$ is nonlocal on the letters $A|_V$, $\lambda$, $c$.
\end{lemma}

\begin{proof}
Compute $Q$ of the tree flow of $\Xi_x$ in two ways. First, $\partial_{\mathsf L_x}$ anticommutes with $Q$ (Lemma \ref{lem:c9-ops}), so by Proposition \ref{prop:c4-defectflow} at order $\hbar^0$ and the classical master equation, $Q$ of $\tfrac\hbar2\partial_{\mathsf L_x}W(P,I)_{\rm tree}$ is $-\tfrac\hbar2\partial_{\mathsf L_x}$ of the tree flow of $\tfrac12\{I,I\}^\sigma_\varepsilon$, the composite closures of the trees of $I$ with one special edge (Proposition \ref{prop:c11-master} with $M=0$). Sort them by the position of the special edge: on the cycle closed by $\mathsf L_x$, which gives the tree flow of $\Xi^\sigma_x$, or on a tree, between two vertices of $I$ or between $\Xi_x$ and a vertex of $I$. Second, with $\Xi_x$ as a vertex, Proposition \ref{prop:c4-defectflow} applied to $I+\delta\,\Xi_x$, $\delta$ of degree $+1$, gives the tree flow of $s\Xi_x$ plus the tree flows with the special-edge parts of $\{I,\Xi_x\}_\varepsilon$ and $\tfrac12\{I,I\}_\varepsilon$ on tree edges. The graphs with the special edge on a tree are the same in both evaluations with the same signs, the sign rule of Proposition \ref{prop:c11-master} being the Koszul rule of the second. Hence the tree flows of $s\Xi_x$ and $\Xi^\sigma_x$ agree, and the tree flow is injective on functionals graded by the number of legs (Theorem \ref{thm:c8-L} (i)). (i) $\Xi^\sigma_x$ is a sum of cycles with one special edge and the composite line, $\sum_m\varepsilon^{-m}a_m+a_0+o(1)$ without $\log\varepsilon$ (Theorem \ref{thm:c11-simplex}); $I^{CT,(1)}(x)$ has no pole (Proposition \ref{prop:c12-poles}), and the coefficient of $\log\varepsilon$ in $s\Xi_x=s\,I^{CT,(1)}(x)+s\Xi^{\rm ren}_x+o(1)$ gives $sL^{(1)}(x)=0$, while the poles of $\Xi^\sigma_x$ vanish. (ii) The coefficient of $\varepsilon^0$.
\end{proof}

\begin{lemma}[Special edges on trees]
\label{lem:c12-BD}
Let $O$ be a local functional. (B) The composite graphs of type I with the special edge on a tree, at an output of the cycle or at a spectator letter of $O$, together with the special-edge parts of the counterterm trees of $\mathrm{ct}^{(1)}_x(O)$, are $e^{-\varepsilon|p|^2}-1$ at a generic momentum $p$ times renormalised composite cycles, and tend to $0$. (D) The composite graphs of type II with the special edge on a tree, and the special-edge parts of the counterterm trees of $I^{CT,(1)}(x)$, tend to $0$.
\end{lemma}

\begin{proof}
(B) Lemma \ref{lem:c12-bridge} with $G$ the composite cycle through the pair with its trees, renormalised by $-\mathrm{ct}^{(1)}_x(O)[\varepsilon]$ joined by the same special edge; at a spectator letter the same holds with the special-edge parts of $\{I,O\}_\varepsilon$ and $\{I,\mathrm{ct}^{(1)}_x(O)\}_\varepsilon$. (D) The graphs are $(e^{-\varepsilon|p|^2}-1)$ times $\Xi_x(\varepsilon)$, which is $\log\varepsilon\,L^{(1)}(x)$ plus a convergent functional (Lemma \ref{lem:c12-Cid} (i)), and $(e^{-\varepsilon|p|^2}-1)\log\varepsilon\to0$.
\end{proof}

\begin{proposition}[The image of the composite cycles through the pairs]
\label{prop:c12-typeI}
For every local functional $O$ and every $\varepsilon>0$,
\[
Q(\text{type I composite graphs of }O)=(\mathrm A)+(\mathrm B)-\epsilon_\Xi\{\Xi_x(\varepsilon),O\}-(\text{composite graphs of }sO)\mp\mathrm{tr}(\Xi_x,sO),
\]
where (A) are the special-edge composite cycles of the pairs of letters of $O$ with trees of $I$ on their outputs (Definition \ref{def:c11-sp}), (B) the type I graphs with the special edge on a tree, $\epsilon_\Xi\in\{\pm1\}$ the sign with which $\{\Xi_x,O\}$ enters $Q\,\mathrm{tr}(\Xi_x,O)$, and the last two terms are present for $sO\neq0$; moreover
\[
Q\,\mathrm{tr}(\Xi_x,O)=\mathrm{tr}(s\Xi_x,O)+\epsilon_\Xi\{\Xi_x,O\}\pm\mathrm{tr}(\Xi_x,sO)+(\mathrm D),
\]
with (D) the type II graphs with the special edge on a tree.
\end{proposition}

\begin{proof}
By Proposition \ref{prop:c11-master} the image under $Q$ of all composite graphs of $O$ is $-\tfrac\hbar2\partial_{\mathsf L_x}$ of the trees with one vertex $sO$ or one special edge; the trees with one vertex $sO$ closed by $\mathsf L_x$ are the composite graphs of $sO$ of both types. The special edge lies on the cycle through two letters (A), on a tree of a type I graph (B), on the cycle of a type II loop (C), or on a tree of a type II graph (D), and (C) is $\mathrm{tr}(\Xi^\sigma_x,O)$. The image of type II, computed with $\Xi_x$ as a vertex by Proposition \ref{prop:c4-defectflow} applied to $I+\delta_1\Xi_x+\delta_2O$, $\delta_1\delta_2$-part, consists of the vertex $s\Xi_x$ with $O$, the vertex $sO$ with $\Xi_x$, the contraction $\{\Xi_x,O\}$ and the special edges on trees, which are (D). Subtract: by Lemma \ref{lem:c12-Cid}, $\mathrm{tr}(s\Xi_x,O)=\mathrm{tr}(\Xi^\sigma_x,O)=$ (C), so (C) and (D) cancel, and $\pm\mathrm{tr}(\Xi_x,sO)$ moves to the right side with the opposite sign.
\end{proof}

\begin{theorem}[The composite term is two-body and local]
\label{thm:c12-local}
Let $O=\int\omega\,X_1\cdots X_n$ be a monomial local functional and $x\in\{\Delta,P_b,M_\omega\}$.
\begin{itemize}
\item (i) The special-edge composite cycles of every pair are $\sum_{m\ge1}\varepsilon^{-m}$ times local coefficients, plus $\mathcal C_x(X_a,X_b)$, plus a remainder tending to $0$; they contain no $\log\varepsilon$.
\item (ii) $I^{CT,(1)}(x)=\log\varepsilon\,L^{(1)}(x)$ with $sL^{(1)}(x)=0$, and the composite part along $x$ is $s\,\Xi^{\rm ren}_x=[\varepsilon^0]\Xi^\sigma_x$.
\item (iii) (Arity-one BRST covariance.) If $sO=0$: $s\,\mathrm{lc}^{(1)}_x(O)=-\epsilon_\Xi\{L^{(1)}(x),O\}$, and $\sum_{a<b}\pm[\varepsilon^{-m}](\text{special-edge composite cycles of }(X_a,X_b))=s\,\mathrm{ct}^{(1)}_{x,m}(O)$ for $m\ge1$.
\item (iv) If $sO=0$, the composite term of $O$ along $x$ exists and
\begin{equation}
(\text{composite term of }O\text{ along }x)=\sum_{a<b}\pm\mathcal C_x(X_a,X_b),
\label{eq:c12-loc}
\end{equation}
a local two-body functional with no one-body term.
\item (v) (Cancellation of the nonlocal one-body terms.) If $sO=0$, the image under $Q$ of the renormalised type I graphs contains $-\epsilon_\Xi\{\Xi^{\rm ren}_x,O\}$ and that of the renormalised type II graphs $+\epsilon_\Xi\{\Xi^{\rm ren}_x,O\}$; both are nonlocal one-body terms on the letters $A|_V$, $\lambda$, $c$ (the factors $\log|p|^2$ and, on $A|_V$, $(p,a)(p,b)|p|^{-2}$ of Remark \ref{rem:c9-Pi}) and local on $A|_{W_{\mathbb R}^\perp}$, and they cancel in the composite term. In the form of Theorem \ref{thm:c10-D1}, the composite terms of the pairs plus those of the legs equal $\sum_{a<b}\pm\mathcal C_x(X_a,X_b)$.
\item (vi) (Gauge-variant operators.) For $sO\neq0$ the right side of \eqref{eq:c12-loc} acquires minus the limit of the renormalised composite graphs of $sO$ of both types minus their counterterm trees: the nonlocal brackets $\{\Xi^{\rm ren}_x,-\}$ on the letters of $sO$ and the $Q$-preimages of the composite cycles through the pairs of letters of $sO$. The components of $O$ are local and two-body by Theorem \ref{thm:c12-master}.
\end{itemize}
\end{theorem}

\begin{proof}
(i) Theorem \ref{thm:c11-simplex} for each chain of Definition \ref{def:c11-sp}. (ii) Lemma \ref{lem:c12-Cid}.

(iii), (iv) The composite term is the limit of $Q$ of the composite graphs minus their counterterm trees, minus the tree insertion of the composite part. Write the counterterm trees as $-\mathrm{tr}(\mathrm{ct}^{(1)}_x(O)[\varepsilon])$ and $-\mathrm{tr}(I^{CT,(1)}(x),O)$. By Proposition \ref{prop:c4-defectflow} and Lemma \ref{lem:c6-bvexact}, with $sO=0$,
\[
\begin{aligned}
Q\big(-\mathrm{tr}(\mathrm{ct}^{(1)}_x(O)[\varepsilon])\big)&=-\mathrm{tr}\big(s\,\mathrm{ct}^{(1)}_x(O)[\varepsilon]\big)+(\text{special edges}),\\
Q\big(-\mathrm{tr}(I^{CT,(1)}(x),O)\big)&=-\mathrm{tr}\big(s\,I^{CT,(1)}(x),O\big)-\epsilon_\Xi\{I^{CT,(1)}(x),O\}+(\text{special edges}),
\end{aligned}
\]
the sign of the bracket being that for $\Xi_x$, both entering as vertices of degree $-1$ joined to $O$ by a tree. By Proposition \ref{prop:c12-typeI} for both types the sum is
\[
(\mathrm A)+(\mathrm B)+(\mathrm D)+(\text{special edges})+\mathrm{tr}\big(s(\Xi_x-I^{CT,(1)}(x)),O\big)-\epsilon_\Xi\{I^{CT,(1)}(x),O\}-\mathrm{tr}\big(s\,\mathrm{ct}^{(1)}_x(O)[\varepsilon]\big),
\]
the brackets $\pm\epsilon_\Xi\{\Xi_x,O\}$ of the two types cancelling. By Lemma \ref{lem:c12-BD}, (B), (D) and the special edges tend to $0$. By Lemma \ref{lem:c12-Cid}, $\mathrm{tr}(s(\Xi_x-I^{CT,(1)}(x)),O)$ tends to the tree insertion of the composite part, which is subtracted in the definition. What remains is (A) $-\epsilon_\Xi\log\varepsilon\{L^{(1)}(x),O\}-\mathrm{tr}(s\,\mathrm{ct}^{(1)}_x(O)[\varepsilon])$. The left side has a limit and (A) has no $\log\varepsilon$ by (i): the coefficient of $\log\varepsilon$ gives $\mathrm{tr}(s\,\mathrm{lc}^{(1)}_x(O)+\epsilon_\Xi\{L^{(1)}(x),O\})=0$, and the coefficient of $\varepsilon^{-m}$ gives $\mathrm{tr}\big(\sum_{a<b}\pm a_m(X_a,X_b)-s\,\mathrm{ct}^{(1)}_{x,m}(O)\big)=0$; the trees with these local vertices are the same at every $\varepsilon$, so injectivity of the tree flow gives (iii). At order $\varepsilon^0$ the trees on the outputs of (A) converge to trees with edges $P(0,T)$, the limit is the tree flow of $\sum_{a<b}\pm\mathcal C_x(X_a,X_b)$, and Theorem \ref{thm:c8-L} (i) gives \eqref{eq:c12-loc}.

(v) Proposition \ref{prop:c12-typeI}: type I contains $-\epsilon_\Xi\{\Xi_x,O\}$, type II $+\epsilon_\Xi\{\Xi_x,O\}$; the counterterm trees move $\Xi_x$ to $\Xi^{\rm ren}_x$ in the second and add $-\epsilon_\Xi\log\varepsilon\{L^{(1)}(x),O\}$, which (iii) removes from the first. The renormalised composite loops on a leg $X$ give the bracket of $\Xi^{\rm ren}_x$ with $X$, whose quadratic part along $\Delta$ is the functional of Remark \ref{rem:c9-Pi}, local on $A|_{W_{\mathbb R}^\perp}$ and with $\log|p|^2$ on $A|_V$, $\lambda$, $c$. The sum of the images of the two types has no bracket, which is (iv); the form of Theorem \ref{thm:c10-D1} splits the same sum into the images of the renormalised composite cycles through the pairs and the brackets $\epsilon_\Xi\{\Xi^{\rm ren}_x,X_a\}$ plus the tree insertion of $s\,\Xi^{\rm ren}_x$.

(vi) For $sO\neq0$, Proposition \ref{prop:c11-master} keeps the trees with one vertex $sO$, whose closures by $\mathsf L_x$ are the composite graphs of $sO$, and the formula for type II keeps $\pm\mathrm{tr}(\Xi_x,sO)$; with their counterterm trees they converge to the renormalised composite graphs of $sO$.
\end{proof}

\begin{remark}[Rotations, and the generators with a non-closed line]
\label{rem:c12-localM}
For $M_\omega$ the composite line is the term $\int(\omega x\cdot\partial A|_V+\omega A|_V,A^*|_V)$ joined to two edges, $Q$-closed because $M_\omega.$ commutes with $Q^{GF}$ and $H$; every step above applies, and by Theorem \ref{thm:c12-pos} (i) $\mathcal C_{M_\omega}=(\omega x_O)^\mu\mathcal C_{P_{e_\mu}}+\mathcal C^0_{M_\omega}$, the intrinsic part computed with the spin term $\tfrac12\omega_{\mu\nu}\Sigma^{\mu\nu}$ and the transport $\omega z\cdot\partial$ centred at $x_O$. Along $K_b$, $\xi$ and $\eta$ the composite line has a nonzero image under $Q$ (Proposition \ref{prop:c12-composite}), Proposition \ref{prop:c11-master} acquires the terms of that image, and the components along these generators are local and two-body by Theorem \ref{thm:c12-master}. The proof of Theorem \ref{thm:c12-local} uses $sO=0$ exactly (or keeps $sO$, (vi)), the classical master equation, $[x.,Q]=0$ with $[Q,P\star x.P\star]=0$, the injectivity of the tree flow on functionals graded by legs, the absence of $\log\varepsilon$ on cycles with one special edge and the weight count that excludes poles of $I^{CT,(1)}(x)$; the sign $\epsilon_\Xi$ enters through the cancellation of the two brackets and is never evaluated.
\end{remark}

